\section{Related work}
\subsection{Prospective scoring and conformal validity}
Dawid's prequential principle evaluates a forecast through the subsequently observed outcome, with the prediction fixed before that outcome is incorporated~\cite{dawid1984}. It specifies an information order rather than a particular fitting algorithm. A repeated full-prefix fit and a native incremental state update can both be evaluated prequentially, although their computational costs differ. This study uses the former. Accordingly, an accurate description of its contribution concerns interval evaluation and numerical reconstruction; it cannot establish that existing forecasting tools lack prospective evaluation or that the selected bases are native stateful models.

Conformal prediction provides prediction sets whose validity is tied to an explicit probabilistic or online model. Vovk, Gammerman and Shafer develop the foundational framework~\cite{vovk2005}; Shafer and Vovk explain both its online interpretation and assumptions~\cite{shafer2008}. An empirical residual quantile resembles one ingredient of that framework, but the resemblance is insufficient to transfer a theorem. A rolling score buffer, repeatedly refitted point model, clipped nominal level and interpolated quantile must be matched to the assumptions and exact construction of a proposed result. The present replay documents those choices rather than assigning exchangeable-data validity to a serially dependent experiment.

Lei and colleagues analyze full and split conformal regression, including efficiency and locally varying interval length~\cite{lei2018}. Their distinction between fitting and calibration information helps explain why in-sample Gaussian residual dispersion is not an independent calibration sample. Romano, Patterson and Cand\`es combine quantile regression with conformal calibration~\cite{cqr2019}. Their construction adapts interval shape through estimated conditional quantiles; the replay instead uses symmetric intervals around a fixed point base and rescales them using past normalized errors. These are different procedures. No conformalized quantile-regression comparator is run here, so its published results are context rather than a numerical baseline in the tables.

Marginal coverage also differs from conditional coverage. Barber and colleagues formalize limitations of unrestricted distribution-free conditional inference~\cite{limits2021}. That result cautions against interpreting a marginal empirical frequency as validity at every time, state or covariate value. It does not excuse avoidable empirical undercoverage or imply that useful conditional improvements are impossible with additional assumptions. The present rolling statistic describes local finite-stream behavior. It is neither a test of every possible conditioning event nor an estimator with an automatically valid confidence band under arbitrary dependence.

\subsection{Dependence and changing distributions}
Several established methods explicitly address departures from exchangeability. Chernozhukov, W\"uthrich and Zhu use structured randomization to accommodate dependent observations, with approximate validity under stated conditions when exchangeability fails~\cite{dependent2018}. Tibshirani and colleagues use weighted conformal prediction for covariate shift, with a known or estimated likelihood ratio~\cite{weighted2019}. Barber and colleagues further study robustness beyond exchangeability and nonsymmetric fitting procedures~\cite{beyond2023}. These approaches introduce assumptions or constructions beyond merely keeping recent scores. The present study estimates no density ratio and runs no block-randomization procedure. Its finite-memory controller should therefore not be presented as an interchangeable implementation of those methods.

Xu and Xie's EnbPI combines bootstrap ensemble predictions and sequential residual information for dynamic time series, with approximate coverage results under the paper's conditions~\cite{enbpi2021}. Their later sequential predictive conformal inference predicts a future residual quantile by exploiting temporal score dependence~\cite{spci2023}. The replay shares an interest in sequential residuals but does not fit an ensemble or a conditional residual-quantile model. Its buffer takes an unconditional interpolated quantile at an adaptively chosen level. EnbPI and SPCI are consequently relevant future comparators, and their omission prevents a claim of superiority over established time-series conformal procedures.

Gibbs and Cand\`es introduce ACI to adapt the effective miscoverage parameter under distribution shift~\cite{gibbs2021}. The replay uses that feedback direction: a miss lowers the effective miscoverage, which usually raises a buffered residual quantile and widens the next interval. However, the replay clips this parameter and uses a particular finite buffer and numerical quantile convention. Those choices define an ACI-style variant rather than proving identity with every setting analyzed in the original paper. The derivation below exposes the additional clipping term in its cumulative-error equation.

Zaffran and colleagues analyze learning-rate effects and propose AgACI through online aggregation~\cite{agaci2022}. Gibbs and Cand\`es subsequently study online conformal prediction with arbitrary distribution shifts and adaptive step-size selection~\cite{gibbs2024}. Both motivate treating gain choice as a substantive design component. The replay retains inherited fixed gains, and its multiplier grid is a sensitivity analysis. It neither implements their adaptive aggregation nor selects a gain on independent validation data. Consequently, reporting the best cell of this grid as a validated improvement would answer a different and unsupported question.

\subsection{Local adaptation, control and forecast evaluation}
Bhatnagar and colleagues develop strongly adaptive online conformal methods whose regret analysis concerns changing local intervals rather than only one whole-horizon average~\cite{strong2023}. This distinction is directly relevant to the difference between averaging signed coverage deviations and averaging their absolute values. Nevertheless, the study's rolling absolute error is a descriptive statistic with a twenty-five-arrival window, not the strongly adaptive regret objective or a reproduction of that paper's algorithms. A lower displayed rolling error cannot establish the corresponding theoretical guarantee.

Angelopoulos, Barber and Bates analyze online conformal prediction with decaying step sizes, distinguishing retrospective coverage from convergence toward a population quantile in stable settings~\cite{decay2024}. The replay has fixed gains and does not estimate the step-size schedule of that method. Angelopoulos, Cand\`es and Tibshirani's Conformal PID Control combines online calibration with prospective score forecasting~\cite{angelopoulos2023}. The historical code considered here instead applies proportional, accumulated and difference terms to a clipped effective miscoverage level. Its PID-style name records that implementation pattern, not a faithful evaluation of the published algorithm. Neither publication's guarantees are imported by naming the controller.

Forecast evaluation must also account for concentration. Gneiting, Balabdaoui and Raftery distinguish calibration from sharpness~\cite{sharpness2007}, and Gneiting and Raftery develop proper scoring-rule principles~\cite{gneiting2007}. The interval score used here penalizes both width and excluded outcomes. It can reveal that a coverage improvement is accompanied by little improvement in the interval's overall score. This is particularly relevant when every controller retains the same point forecast. The study consequently reports scaled interval score and width rather than presenting coverage as a complete measure of probabilistic forecasting quality.

Finally, forecasting evaluation depends on origin choice, information availability and aggregation. Tashman's analysis of out-of-sample evaluation motivates retaining temporal order and explicit evaluation schedules~\cite{tashman2000}. Hyndman and Koehler's discussion of error measures motivates scrutiny of scaling and zero denominators~\cite{mase2006}. The replay uses a warmup-only seasonal scale for cross-case interval comparisons and retains the historical per-arrival quantile-loss ratio only for archive reconciliation. These decisions prevent a later outcome from setting the scale and expose the unstable behavior of a denominator near zero. They do not establish that the selected fifteen streams represent an operational population or that their average score identifies a generally optimal interval controller.
